\part*{Preámbulo}

Tomé estos apuntes cuando cursé Introducción a la Programación en la Universidad de los Andes, entre agosto y noviembre de 2020. Fueron refinados en mi tiempo como monitor en 2021, luego como tutor en el centro de enseñanza de programación de la universidad en 2022, y tras eso como asistente de investigación en ese mismo centro desde el 2022 hasta el 2024. Finalmente, fueron revisados en mi etapa como profesor de la asignatura entre el 2025 y 2026, pero no fueron modificados sustancialmente\footnote{De hecho, los apuntes están presentados en el mismo formato en el que originalmente fueron escritos; la clase de \LaTeX{} que determina su estilo la escribí en noviembre del 2020.}: siguen siendo apuntes de un estudiante para estudiantes. 

\section{Programar}

Programar es, en esencia, escribir instrucciones que puedan ser ejecutadas por una máquina.

En general, escribir instrucciones para una máquina es más difícil que escribirlas para un humano por (al menos) tres razones:
\begin{itemize}
  \item La máquina no entiende español, por lo que las instrucciones deben estar escritas en un lenguaje diseñado para que la máquina entienda. Existen centenas de dichos \emph{lenguajes de programación}. \emph{Python} es uno de ellos.
  \item Los lenguajes naturales (como el español) tienen reglas que determinan cómo organizar las palabras para que esté correctamente escrito. A eso se le llama \emph{sintaxis}. Es fácil para un humano entender español ligeramente mal escrito, pero una máquina no puede interpretar Python con sintaxis incorrecta.
  \item El significado de un conjunto de palabras es su \emph{semántica}. Los humanos pueden entender la intención detrás de instrucciones mal formuladas, pero una máquina ejecuta lo escrito de forma literal, así tenga poco sentido o no haga lo que se pretendía.
\end{itemize}

Esos tres factores hacen que programar sea difícil y que sea ampliamente considerado un arte.
%TODO: Hyperlink a que la programación es un arte.
% TODO: Concluir mejor esta sección. Es difícil pero y qué

\begin{tip}
  En estos apuntes se usa mucho la palabra sintaxis. En esencia, es el orden en que se escriben las cosas, que es muy importante al programar.
\end{tip}
